\chapter{Diferencimi dhe integrimi numerik}
\begin{qellime}
Studenti nxjerr formulat e diferencave të fundme; analizon gabimin e këputjes dhe rrumbullakimit; zbaton trapezin, Simpsonin dhe kuadraturën e Gausit; zgjedh metodën sipas lëmueshmërisë dhe dimensionit.
\end{qellime}

\section{Derivatet nga zhvillimi Taylor}
Nga
\begin{equation}
 f(x+h)=f(x)+hf'(x)+\frac{h^2}{2}f''(x)+\cdots
\end{equation}
merret diferenca përpara
\begin{equation}
 f'(x)=\frac{f(x+h)-f(x)}{h}+\Oh(h),
\end{equation}
ndërsa kombinimi simetrik jep
\begin{equation}
 f'(x)=\frac{f(x+h)-f(x-h)}{2h}+\Oh(h^2).
\end{equation}
Për derivatin e dytë,
\begin{equation}
 f''(x)=\frac{f(x+h)-2f(x)+f(x-h)}{h^2}+\Oh(h^2).
\end{equation}
Koeficientët mund të ndërtohen sistematikisht duke kërkuar që kombinimi linear të riprodhojë polinomet deri në një gradë të caktuar.

\section{Hapi optimal}
Gabimi i këputjes ulet me $h$, por zbritja $f(x+h)-f(x)$ humbet shifra kur $h$ është shumë i vogël. Për diferencën përpara,
\begin{equation}
 E(h)\approx C_1h+C_2\frac{u}{h},
\end{equation}
prandaj ekziston një hap optimal i rendit $\sqrt{u}$. Për diferencën qendrore, kompromisi tipik jep $h\sim u^{1/3}$.

\begin{figure}[ht]
\centering\includegraphics[width=.74\textwidth]{08_gabimi_derivimit.pdf}
\caption{Gabimi total ka formë U: këputja dominon për hapa të mëdhenj, rrumbullakimi për hapa tepër të vegjël.}
\end{figure}

\section{Kuadratura e përbërë}
Integrali ndahet në nënintervale. Rregulla e trapezit është
\begin{equation}
 T_n=h\left[\frac{f(a)+f(b)}2+\sum_{i=1}^{n-1}f(a+ih)\right],\qquad h=\frac{b-a}{n},
\end{equation}
me gabim global $\Oh(h^2)$. Rregulla e Simpsonit për $n$ çift është
\begin{equation}
 S_n=\frac h3\left[f_0+f_n+4\sum_{i\text{ tek}}f_i+2\sum_{i\text{ çift},\,i\ne0,n}f_i\right],
\end{equation}
me gabim global $\Oh(h^4)$ për funksione mjaft të lëmuara.

\begin{figure}[ht]
\centering\includegraphics[width=.74\textwidth]{09_kuadratura_konvergjenca.pdf}
\caption{Konvergjenca e trapezit dhe Simpsonit për një integral me zgjidhje të njohur.}
\end{figure}

\section{Kuadratura e Gausit}
Rregulli me $n$ nyje
\begin{equation}
 \int_{-1}^{1}f(x)\dd x\approx\sum_{i=1}^{n}w_if(x_i)
\end{equation}
zgjedh nyjet si zerot e polinomit Legendre $P_n$. Ai integron saktë çdo polinom me gradë deri $2n-1$. Një transformim afin e çon intervalin $[a,b]$ në $[-1,1]$. Kuadratura Gauss është e fuqishme kur funksioni është i lëmuar dhe dimensioni i vogël.

\section{Integrale të papërshtatshme dhe shumëfishta}
Singularitetet e integrueshme trajtohen me ndryshim ndryshoreje ose ndarje të intervalit. Intervali i pafund mund të hartohet në interval të fundmë. Në $d$ përmasa, prodhimi tensorial i $n$ nyjeve kërkon $n^d$ vlerësime: mallkimi i dimensionalitetit. Kjo është pika ku metodat Monte Carlo bëhen konkurruese, sepse gabimi i tyre bazë varet nga numri i mostrave dhe jo drejtpërdrejt nga $d$.

\begin{shembull}[Perioda e lavjerrësit]
Për amplitudë $\theta_0$, perioda e saktë e modelit ideal është
\begin{equation}
 T=4\sqrt{\frac Lg}\int_0^{\pi/2}\frac{\dd\phi}{\sqrt{1-k^2\sin^2\phi}},\qquad k=\sin(\theta_0/2).
\end{equation}
Integrali eliptik mund të llogaritet me Simpson ose Gauss dhe të krahasohet me $T_0=2\pi\sqrt{L/g}$. Diferenca mat kufirin e përafrimit të këndit të vogël.
\end{shembull}

\begin{lidhje}
Laboratori llogarit periodën për disa amplituda, kryen studim konvergjence dhe ndan gabimin e kuadraturës nga gabimi i modelit $\sin\theta\approx\theta$.
\end{lidhje}

\section*{Pyetje kontrolli}
\begin{enumerate}
\item Pse derivati numerik nuk përmirësohet pafundësisht kur $h\to0$?
\item Çfarë kushti duhet për rendin e katërt të Simpsonit?
\item Pse prodhimi tensorial bëhet i papërshtatshëm në dimension të lartë?
\end{enumerate}
\section{Algoritmet e diferencimit dhe kuadraturës}
\subsection{Derivati qendror dhe studimi i hapit}
\begin{lstlisting}[language=Python,caption={Gabimi i derivatit qendror për një varg hapash.}]
import numpy as np

def derivat_qendror(f, x, h):
    return (f(x + h) - f(x - h)) / (2.0*h)

f = np.sin
x0 = 1.0
h = np.logspace(-16, -1, 200)
D = derivat_qendror(f, x0, h)
gabimi = np.abs(D - np.cos(x0))
h_opt = h[np.argmin(gabimi)]
print("h optimal empirik =", h_opt)
\end{lstlisting}
Grafiku $E(h)$ duhet të paraqitet në bosht log--log. Në zonën e këputjes pritet pjerrësi afërsisht 2; në zonën e rrumbullakimit gabimi rritet kur $h$ zvogëlohet.

\subsection{Rregulli i përbërë i Simpsonit}
\begin{lstlisting}[language=Python,caption={Integrimi Simpson me kontrollin e numrit çift të nënintervaleve.}]
def simpson(f, a, b, n):
    if n <= 0 or n % 2 != 0:
        raise ValueError("n duhet te jete numer pozitiv çift.")
    x = np.linspace(a, b, n + 1)
    y = f(x)
    h = (b - a) / n
    return h/3 * (y[0] + y[-1]
                  + 4*np.sum(y[1:-1:2])
                  + 2*np.sum(y[2:-2:2]))

for n in [4, 8, 16, 32, 64]:
    I = simpson(np.sin, 0.0, np.pi, n)
    print(n, I, abs(I - 2.0))
\end{lstlisting}

\subsection{Vlerësimi Runge i gabimit}
Nëse metoda ka rend $p$, dy llogaritje me $h$ dhe $h/2$ japin
\begin{equation}
 E_{h/2}\approx\frac{|A(h/2)-A(h)|}{2^p-1}.
\end{equation}
Ky vlerësim është i vlefshëm vetëm në regjimin asimptotik. Prandaj kontrollohet edhe raporti i gabimeve për tri hapa radhazi.

\begin{lstlisting}[language=Python]
def studim_konvergjence(llogarit, ns, p):
    v = np.array([llogarit(n) for n in ns])
    e_runge = np.abs(v[1:] - v[:-1]) / (2**p - 1)
    return v, e_runge
\end{lstlisting}

\subsection{Kuadratura e Gausit}
\begin{lstlisting}[language=Python,caption={Gauss--Legendre në një interval të përgjithshëm.}]
from numpy.polynomial.legendre import leggauss

def gauss_legendre(f, a, b, n):
    xi, w = leggauss(n)
    x = 0.5*(b-a)*xi + 0.5*(a+b)
    return 0.5*(b-a) * np.sum(w * f(x))
\end{lstlisting}
Në notebook krahasohen numri i vlerësimeve të funksionit, gabimi dhe koha e ekzekutimit. Kjo e lidh kuadraturën klasike me motivimin për Monte Carlo në dimension të lartë.

Formulat dhe emërtimet shqip të diferencimit e integrimit numerik mbështeten te \cite{hanelli,hoxha,datja}; interpretimi fizik dhe krahasimi me Monte Carlo plotësohen nga \cite{newman,press}.

\subsection{Nxjerrja sistematike e formulave të diferencimit}
Nga zhvillimet Taylor
\begin{align}
 f(x+h)&=f(x)+hf'(x)+\frac{h^2}{2}f''(x)+\frac{h^3}{6}f'''(x)+\Oh(h^4),\\
 f(x-h)&=f(x)-hf'(x)+\frac{h^2}{2}f''(x)-\frac{h^3}{6}f'''(x)+\Oh(h^4)
\end{align}
dalin dy kombinime të rëndësishme. Zbritja e barazimeve anulon termat çift dhe jep
\begin{equation}
 f'(x)=\frac{f(x+h)-f(x-h)}{2h}-\frac{h^2}{6}f'''(x)+\Oh(h^4),
\end{equation}
pra diferenca qendrore ka gabim të këputjes $\Oh(h^2)$. Mbledhja anulon termat tek dhe jep
\begin{equation}
 f''(x)=\frac{f(x+h)-2f(x)+f(x-h)}{h^2}-\frac{h^2}{12}f^{(4)}(x)+\Oh(h^4).
\end{equation}
Koeficientët e formulave me më shumë pika gjenden duke kërkuar që kombinimi linear të riprodhojë saktë monomet $1,x,x^2,\ldots$ deri në gradën e dëshiruar.

Gabimi i rrumbullakimit në diferencën e parë rritet afërsisht si $u|f|/h$, ndërsa gabimi i këputjes për diferencën qendrore ulet si $C h^2$. Modeli
\begin{equation}
 E(h)\simeq C_1h^2+C_2\frac{u}{h}
\end{equation}
ka minimum për $h_{\mathrm{opt}}\sim u^{1/3}$. Kjo shpjegon formën në U të grafikut të gabimit: zvogëlimi i pakufizuar i hapit nuk rrit saktësinë.

\subsection{Rregulla e trapezit dhe termi i gabimit}
Në intervalin $[a,b]$, polinomi linear interpolues është
\begin{equation}
 p_1(x)=f(a)\frac{b-x}{b-a}+f(b)\frac{x-a}{b-a}.
\end{equation}
Integrimi i tij jep
\begin{equation}
 T_1=\frac{b-a}{2}[f(a)+f(b)].
\end{equation}
Nga mbetja e interpolimit ose nga Taylor-i rreth mesit merret
\begin{equation}
 \int_a^b f(x)\dd x-T_1=-\frac{(b-a)^3}{12}f''(\xi).
\end{equation}
Për $n$ nënintervale me hap $h=(b-a)/n$, mbledhja e gabimeve lokale jep gabim global
\begin{equation}
 E_T=-\frac{b-a}{12}h^2 f''(\xi),
\end{equation}
pra rregulla e përbërë është e rendit të dytë.

\subsection{Rregulla e Simpsonit}
Mbi dy nënintervale, integrohet polinomi kuadratik që kalon në $x_0$, $x_1=x_0+h$ dhe $x_2=x_0+2h$:
\begin{equation}
 S_2=\frac{h}{3}[f(x_0)+4f(x_1)+f(x_2)].
\end{equation}
Edhe pse interpoluesi ka gradë dy, simetria e bën rregullën të saktë edhe për polinomet kubike. Gabimi lokal është
\begin{equation}
 -\frac{h^5}{90}f^{(4)}(\xi),
\end{equation}
ndërsa përbërja në një interval të fiksuar jep gabim global $\Oh(h^4)$. Numri i nënintervaleve duhet të jetë çift. Një kod i mirë e verifikon këtë parakusht, në vend që të prodhojë në heshtje një rezultat të pakuptimtë.

\subsection{Kuadratura e Gausit}
Kërkohen nyje $x_i$ dhe pesha $w_i$ të tilla që
\begin{equation}
 \int_{-1}^{1}f(x)\dd x\approx\sum_{i=1}^{n}w_if(x_i)
\end{equation}
të jetë e saktë për gradën më të lartë të mundshme. Me $2n$ parametra përcaktohen momentet deri në gradën $2n-1$. Nyjet optimale janë rrënjët e polinomit Legendre $P_n$. Për $n=2$, simetria kërkon nyje $\pm a$ dhe pesha të barabarta $w$; saktësia për $1$ dhe $x^2$ jep $2w=2$ dhe $2wa^2=2/3$, prandaj $w=1$ dhe $a=1/\sqrt3$.

Për intervalin $[a,b]$ përdoret transformimi
\begin{equation}
 x=\frac{a+b}{2}+\frac{b-a}{2}\xi,qquad
 \dd x=\frac{b-a}{2}\dd\xi.
\end{equation}
Avantazhi është saktësia e lartë me pak vlerësime; kufizimi është se nyjet nuk janë të futura natyrshëm, prandaj kontrolli adaptiv kërkon çifte rregullash ose nënndarje.

\begin{lstlisting}[language=Python,caption={Simpson i përbërë me kontroll të dyfishimit të rrjetit.}]
import numpy as np

def simpson(f, a, b, n):
    if n <= 0 or n % 2:
        raise ValueError("n duhet të jetë numër pozitiv çift.")
    x = np.linspace(a, b, n + 1)
    h = (b - a) / n
    return (h / 3.0) * (f(x[0]) + f(x[-1])
           + 4.0*np.sum(f(x[1:-1:2]))
           + 2.0*np.sum(f(x[2:-1:2])))

def kontroll_simpson(f, a, b, n):
    I_h = simpson(f, a, b, n)
    I_h2 = simpson(f, a, b, 2*n)
    gabimi = abs(I_h2 - I_h) / (2.0**4 - 1.0)
    return I_h2, gabimi
\end{lstlisting}

\subsection{Përmbledhje dhe ushtrime}
Derivimi numerik është intrinsikisht më i ndjeshëm ndaj zhurmës, ndërsa integrimi ka efekt mesatarizues. Rendi formal duhet konfirmuar me një studim rrjeti dhe jo vetëm cituar.
\begin{enumerate}
\item Nxirrni formulën qendrore me pesë pika për $f'(x)$ dhe termin kryesor të gabimit.
\item Vërtetoni saktësinë kubike të Simpsonit duke integruar monomet $1,x,x^2,x^3$.
\item Ndërtoni Gauss--Legendre me tri pika nga momentet dhe simetria.
\item Llogaritni periodën e lavjerrësit si integral dhe krahasoni Simpsonin me Gauss-in.
\end{enumerate}

\subsection{Kuadratura adaptive dhe vlerësimi lokal i gabimit}
Një rrjet uniform harxhon pika në rajone ku funksioni është i lëmuar dhe mund të mos zgjidhë rajone me ndryshim të shpejtë. Kuadratura adaptive e ndan intervalin aty ku gabimi i vlerësuar është i madh. Për Simpsonin, le të jenë $S(a,b)$ vlerësimi në të gjithë intervalin dhe
\begin{equation}
 S_2=S(a,m)+S(m,b),\qquad m=\frac{a+b}{2}.
\end{equation}
Meqë gabimi kryesor shkallëzohet si $h^4$ globalisht, vlerësimi i Richardson-it është
\begin{equation}
 E\approx\frac{S_2-S(a,b)}{2^4-1}.
\end{equation}
Nëse $|E|$ është nën tolerancën lokale, pranohet $S_2+E$; përndryshe të dy gjysmat ndahen rekursivisht. Toleranca globale ndahet ndërmjet intervaleve, për shembull përgjysmohet në çdo degë. Duhet vendosur thellësi maksimale për të shmangur rekursion të pafundëm pranë singulariteteve ose zhurmës.

Vlerësimi mund të dështojë kur dy rregullat japin rastësisht të njëjtin rezultat, kur funksioni ka një kulm shumë të ngushtë që nuk kapet nga pikat ose kur ka ndërprerje. Një algoritëm profesional kombinon rregulla të futura Gauss--Kronrod, zbulim të sjelljes së keqe dhe raport të statusit, jo vetëm një numër.

\subsubsection{Integralet e papërshtatshme}
Për një kufi të pafundëm, transformimi
\begin{equation}
 x=\frac{t}{1-t},\qquad \dd x=\frac{\dd t}{(1-t)^2},\qquad 0\le t<1,
\end{equation}
e kthen $[0,\infty)$ në $[0,1)$. Transformimi duhet zgjedhur sipas shuarjes së integrandit; Jakobi mund të krijojë një funksion të vështirë pranë $t=1$. Një alternativë është prerja në $R$ dhe kufizimi analitik i bishtit.

Për singularitetin integrueshëm $x^{-1/2}$ në zero, zëvendësimi $x=t^2$ jep
\begin{equation}
 \int_0^1\frac{g(x)}{\sqrt x}\dd x
 =2\int_0^1g(t^2)\dd t,
\end{equation}
që është i rregullt nëse $g$ është i lëmuar. Njohja e strukturës analitike është më efikase se zvogëlimi brutal i hapit.

\subsubsection{Integralet shumëpërmasore}
Rregulli tensorial me $n$ pika për bosht kërkon $n^d$ vlerësime në $d$ përmasa. Kjo është mallkimi i dimensionalitetit. Për dimensione të ulëta dhe integrandë të lëmuar, kuadratura tensoriale ose rrjetet e rralla mund të jenë shumë të sakta. Për dimensione të larta, Monte Carlo me gabim $N^{-1/2}$ bëhet konkurrues sepse eksponenti nuk varet shprehimisht nga $d$.

Megjithatë, dimensioni ndikon përmes variancës dhe gjeometrisë. Nëse rajoni $D$ zë një pjesë jashtëzakonisht të vogël të kutisë kufizuese, metoda hit-or-miss ka variancë të madhe. Një parametrizim që gjeneron pika drejtpërdrejt në $D$ është shumë më efikas. Për një disk, përdorimi i $r=R\sqrt U$ dhe $\theta=2\pi V$ prodhon shpërndarje uniforme në sipërfaqe; zgjedhja $r=RU$ do të grumbullonte pika pranë qendrës.

\subsubsection{Verifikimi me funksione prove}
Një rutinë integrimi testohet me polinome të gradës së saktësisë, funksione periodike, kulme të ngushta dhe raste me zgjidhje të njohur. Rendi matet duke vizatuar gabimin ndaj hapit në shkallë log--log. Pjerrësia duhet të jetë pranë rendit teorik vetëm para se rrumbullakimi të dominojë. Raporti duhet të përmbajë numrin e vlerësimeve të funksionit; dy metoda me të njëjtin gabim mund të kenë kosto shumë të ndryshme.

